% !TeX program = xelatex
% ======================================================================
% 天津中德应用技术大学 软件工程专业 本科毕业设计（论文）
% 使用说明：先填写下面的论文信息，然后依次撰写各部分。
% 编译：latexmk（自动调用 xelatex + bibtex），或手动执行
%   xelatex main && bibtex main && xelatex main && xelatex main
% ======================================================================
\documentclass[fandol]{izdthesis}

% ---------------- 论文信息（封面、页眉使用） ----------------
\ctitle{基于×××的×××系统的设计与实现}        % 中文题目（宋体三号）
\etitle{Design and Implementation of ××× System Based on ×××} % 外文题目
\stuname{张三}                              % 姓名
\college{软件与通信学院}                     % 学院
\major{软件工程}                            % 专业
\advisor{李四}                              % 指导教师
\advisortitle{副教授}                        % 职称
\finishdate{2026 年 5 月}                   % 完成时间
\gradyear{2026}                            % 届别（页眉使用）

\begin{document}

% ---------------- 封面 ----------------
\makecover

% ---------------- 声明页 ----------------
\makedeclaration

% ---------------- 中文摘要（600 字左右，至少三段，一页之内，约占 80% 篇幅） ----------------
\begin{cnabstract}
  在此撰写中文摘要。摘要应高度凝练本毕业设计的工作、特色与背景，
  说明系统解决了什么问题、采用了什么技术路线、取得了什么效果。
  第一段一般介绍研究背景与问题提出。

  第二段一般介绍系统的需求分析、总体设计、关键技术选型与实现过程，
  说明完成的主要功能模块。

  第三段一般介绍系统测试情况与最终成果，说明系统达到的预期目标。

  \keywords{×××××；×××××；×××××；×××××；×××××}
\end{cnabstract}

% ---------------- 英文摘要（与中文摘要对应） ----------------
\begin{enabstract}
  The English abstract should correspond to the Chinese abstract.
  It describes the background, the problem to be solved, the technical
  approach, the main functions of the system, and the final results.

  Write three or more paragraphs here, mirroring the structure of the
  Chinese abstract.

  \enkeywords{×××××；×××××；×××××；×××××；×××××}
\end{enabstract}

% ---------------- 目录 ----------------
\izdtableofcontents

% ---------------- 正文（从此处开始有页眉、页码） ----------------
\izdmainmatter

\chapter{引言}

\section{研究背景}
描述本毕业设计的背景领域内容，提出当前领域的困难，说明本毕业设计
解决了什么问题。文献引用多在此节，引用在正文中按顺序以上标标注，
例如：随着微服务架构的普及\cite{example2024microservice}，……
文中出现的任何数据都要有明确出处。

\section{国内外研究现状}
阐述当前研究领域的现状。

\subsection{国内研究现状}
国内学者在×××方面开展了大量研究\cite{example2024survey}。

\subsection{国外研究现状}
国外研究起步较早\cite{example2023foreign}。

\section{本课题的目的及意义}
阐述当前领域内的问题，以及本毕业设计能够解决什么、提升什么。

\section{本文主要工作}
介绍本毕业设计完成的主要工作与系统功能内容。

\section{组织结构}
介绍本论文的章节安排。

\chapter{相关技术}
本章说明毕业设计用到的技术并做介绍，用自己的语言重新组织复述，
一般 2 页即可。例如：××系统后端采用 Spring Boot
框架\cite{example2024springboot}，如图~\ref{fig:arch} 所示。

\begin{figure}[htbp]
  \centering
  % \includegraphics[width=0.8\textwidth]{figures/example.png}
  \fbox{\rule{0pt}{3cm}\rule{6cm}{0pt}} % 占位框，替换为上面的图片
  \caption{系统整体架构图}
  \label{fig:arch}
\end{figure}

\chapter{需求分析}
本章描述需求阶段的工作，包括功能性需求分析与非功能性需求分析，
必须有功能的用例分析图，系统功能一般在 10 个左右。

\section{功能性需求分析}
针对每个功能进行用户需求分析，每个功能至少一段话的描述。

\section{非功能性需求分析}
从独立性、复杂性、扩展性、容错性等方面分析（注意查重）。

\section{用例分析}
针对功能性需求画用例图，并对每个用例撰写用例描述，如表~\ref{tab:usecase} 所示。

\begin{table}[htbp]
  \centering
  \caption{用户登录用例描述表}
  \label{tab:usecase}
  \zihao{5}\songti
  \begin{tabular}{|l|l|}
    \hline
    用例名称 & 用户登录 \\ \hline
    参与者   & 用户     \\ \hline
    前置条件 & 用户已注册 \\ \hline
    基本流程 & 输入账号密码，系统校验并跳转 \\ \hline
    后置条件 & 登录成功，进入系统首页 \\ \hline
  \end{tabular}
\end{table}

\section{活动图}
针对用例图画出活动图（可选）。

\chapter{系统设计}

\section{系统总体设计}
用包图体现整体技术架构，辅以类图、构件图描述功能实现的相关类。
系统负载均衡的加权评分可按公式~\ref{eq:score} 计算。

\begin{equation}
  S = \alpha \cdot \frac{1}{T} + \beta \cdot C
  \label{eq:score}
\end{equation}

\section{系统详细设计}
分别详细描述每个末级功能，并配有功能时序图（一般 10 个左右）。

\section{数据库设计}
对整体数据库做 E-R 图并加以描述，列出每张数据库表的表结构
（10 张表左右，过多可摘选关键表）。

\chapter{系统实现}
说明项目包结构，上面章节描述的每一个功能都应在本章贴出实现图，
并可粘贴关键代码（不超过半页），配以文字说明其作用，例如：

\begin{lstlisting}[language=Java]
@RestController
@RequestMapping("/api/user")
public class UserController {
    @Autowired
    private UserService userService;

    @PostMapping("/login")
    public Result login(@RequestBody LoginDTO dto) {
        return userService.login(dto);
    }
}
\end{lstlisting}

\chapter{系统测试}
进行系统功能测试，撰写测试用例，可按等价类划分法或边界值分析法。

\chapter{总结与展望}
总结本毕业设计的工作与内容，指出缺陷并提出展望。本章应超过 1 页。

% ---------------- 参考文献（不少于15篇，近5年，外文不少于2篇） ----------------
\izdbibliography{ref}

% ---------------- 致谢（限 1 页） ----------------
\begin{acknowledgements}
  对毕业设计（论文）的指导教师，提过有益建议或给予过帮助的老师、
  同学以及其他人，在论文结尾书面致谢。
\end{acknowledgements}

% ---------------- 附录（可选，较长代码放此处，不计入正文字数） ----------------
% \begin{izdappendix}
%   在此放置较长的源代码或其他补充材料。
% \end{izdappendix}

\end{document}
